\newpage
\chapter*{\centering Abstract}
\addcontentsline{toc}{chapter}{Abstract}

\quad Chain-of-thought (CoT) prompting and search-based reasoning represent two contrasting inference-time paradigms for large language models. CoT commits to a single linear trajectory, while search explores alternative partial reasoning paths before committing to an answer. In this thesis, we study this trade-off using A* as a controlled instantiation of heuristic-guided search.

Across eight benchmarks---StrategyQA, LogicQA, HotpotQA, GSM8K, MATH, HumanEval, MBPP, and CodeContests---spanning question answering, mathematical reasoning, and code generation, and across three model scales (0.5B, 8B, 14B), we find that search is not a universal upgrade over CoT. Instead, the two strategies solve substantially different subsets of instances, with oracle gaps of up to 14 percentage points above the best single method.

To explain this complementarity, we introduce a \emph{Reasoning Topology Framework} based on 13 structural features computable from the input alone, before any generation. A derived branching coefficient predicts when search helps better than dataset identity ($R^2 = 0.94$ on QA/math clusters). Building on this, we train a topology-aware router that selects between CoT and A* without observing generated traces. Despite this restriction, it outperforms post-hoc routers based on semantic entropy and LLM-as-critic judgments, recovering 32.3\% of the oracle gap at 8B scale across all benchmarks.

We attribute this counterintuitive result to a \emph{Fluency--Correctness Asymmetry}: in 38\% of disagreement cases, the more fluent trace is actually wrong, systematically misleading post-hoc judges. On code generation, the topology router recovers 67.8\% of the HumanEval oracle gap and 89.7\% on CodeContests at 0.5B scale. Our results suggest that adaptive reasoning should route by problem structure rather than trace evaluation, deploying search only when the input topology warrants its higher cost.

\vskip 2ex
\noindent\textbf{Keywords:} Large Language Models, Chain-of-Thought, A* Search, Reasoning Topology, Adaptive Routing, Code Generation
